Real-time collaborative applications have become part of everyday life, letting people communicate and work together across distance and time. A widely used example is the document editor that several users can edit at once. This thesis presents PeerToCP, a local-first collaborative code editor built on WebRTC and conflict-free replicated data types (CRDTs), together with a shared shell that one user runs and every other user in the network group can use. Three backend variants of the application are compared. For keeping documents consistent, two approaches are evaluated: operational transformation (OT) and CRDTs. For the network architecture, the study evaluates a client-server CRDT variant, a peer-to-peer CRDT variant and a client-server OT variant. A peer-to-peer OT variant was not evaluated, because the OT implementation used here needs a server as its single source of truth. The peer-to-peer CRDT variant performed best for groups of up to eight users ($n \leq 8$), and its signaling server uses minimal resources, so one server can host more network groups. The client-server CRDT variant, however, grows more slowly with the number of users in a group and uses fewer resources on the clients, at the cost of more on the server. It is therefore worth considering when a peer-to-peer network cannot be established, or when groups are much larger than those tested in this study.

\vspace{1em}
\noindent\textbf{Keywords:} WebRTC, WebSocket, local-first, real-time, code editor, conflict-free replicated data types, operational transformation, shared shell, Yjs, Electron
